\section*{Erd\H{o}s Problem 419}

\paragraph{Statement.}
Determine the limit points of \(R(n)=d((n+1)!)/d(n!)\), where \(d\)
counts positive divisors. They are exactly
\[
 \{1\}\ \cup\ \{1+1/k:k=1,2,\ldots\}.
 \tag{1}
\]
The result is known from Erd\H{o}s, Graham, Ivi\'c and Pomerance.
The following elementary reconstruction follows the prime-factor approach
recorded on the problem page.

\paragraph{Small prime factors contribute a factor tending to one.}
The divisor formula gives
\[
 R(n)=\prod_{p\mid n+1}
 \left(1+\frac{v_p(n+1)}{v_p(n!)+1}\right).
\]
For \(p\leq n^{2/3}\),
\[
 v_p(n!)+1\geq\lfloor n/p\rfloor+1>n/p,\qquad
 v_p(n+1)\leq\log_2(n+1).
\]
There are at most \(\log_2(n+1)\) distinct prime factors. Hence the
product over these small primes lies between \(1\) and
\[
 \exp\left(O\left(\frac{(\log n)^2}{n^{1/3}}\right)\right)=1+o(1),
 \tag{2}
\]
uniformly in \(n\).

\paragraph{There is at most one remaining factor.}
For sufficiently large \(n\), two prime factors exceeding \(n^{2/3}\),
or the square of one such prime, would exceed \(n+1\).
If there is no such prime, (2) gives \(R(n)=1+o(1)\).
Otherwise write \(n+1=kp\), with \(p>n^{2/3}\).
Its exponent in \(n+1\) is one, and \(p^2>n\), so
\[
 v_p(n!)=\lfloor n/p\rfloor=k-1.
\]
The remaining factor is exactly \(1+1/k\). Thus \(R(n)\) has distance
tending to zero from the set in (1). That set is closed; all limit
points of \(R(n)\) belong to it.

\paragraph{Every proposed limit point occurs.}
Fix \(k\geq1\), let \(p\) run through sufficiently large primes, and put
\(n=kp-1\). Then \(p>n^{2/3}\), so (2) yields \(R(n)\to1+1/k\).
For \(1\), take \(n=2^j-1\). Legendre's formula gives
\(v_2(n!)=2^j-1-j\), whence
\[
 R(2^j-1)=1+\frac{j}{2^j-j}\longrightarrow1.
\]
This proves both directions of (1).

\paragraph{Source and scope.}
\url{https://www.erdosproblems.com/419}, checked 24 September 2026,
records the known answer and an exposition attributed to Mehtaab Sawhney.
The proof here uses only prime factorization, Legendre's formula, and
infinitude of primes. No new result or formal verification is claimed.

\paragraph{Completion estimate.}
The complete limit-point classification is proved.
\textbf{COMPLETION: 100\%}
